\documentclass[12pt]{article}
\usepackage[utf8]{inputenc}
\usepackage[T1]{fontenc}
\usepackage{hyperref}
\title{Coahuilteco Language Resources for K-12 Curriculum Development}
\author{Compiled by Claude Sonnet 5.5 (claude-sonnet-5-5) for Andr\'as Kornai}
\date{}
\begin{document}
\maketitle

\section{Introduction}
Coahuilteco (ISO 639-3 xcw; Glottocode coah1252) is an extinct language of northern Mexico and southern Texas. Its classification is disputed. It was a lingua franca among Coahuiltecan peoples at the San Antonio missions. My searches found no school programmes, living-speaker communities or classroom materials.

\section{Primary Educational Resources}
None found. I saw no primers, textbooks or school materials for Coahuilteco.

\section{Linguistic and Reference Materials}
The main reference is Swanton's \emph{Coahuilteco-English Vocabulary} (Bureau of American Ethnology Bulletin 127, 1940), mentioned in the sources I saw. A search result pointed to the Biodiversity Heritage Library as a place to find Bureau of American Ethnology material, but I did not see a direct URL for the vocabulary and did not read its text. Overviews are at \url{https://en.wikipedia.org/wiki/Coahuiltecan_languages}, \url{https://en.wikipedia.org/wiki/Comecrudan_languages} (the Wikipedia page for the related Garza language is at \url{https://en.wikipedia.org/wiki/Garza_language}), \url{https://www.indigenouscultures.org/coahuiltecan-language} and \url{https://asjp.clld.org/languages/COAHUILTECO}.

\section{Religious and Community Resources}
None found in my searches. The language was recorded at the Catholic missions in San Antonio, but I saw no URL for any mission-era text.

\section{Digital and Audio Resources}
No audio recordings were found. The ASJP entry above is a short word list in a comparative database.

\section{Curriculum Development Context}
The documentation is thin, so a K-12 programme would be mostly cultural and historical. It could teach the Coahuiltecan peoples, mission-era history and how languages are documented. Teaching the language itself would require specialist reconstruction from Swanton's vocabulary and consultation with descendant communities. I found no sign that either exists.
\end{document}